% Transcription — fonds Alexandre Grothendieck, shelfmark 36, batch 3 (pages 41-60).
% Established from the facsimile archives/batches/36/batch-03.pdf (a copy of
% 36.pdf fetched through the project's relay, grothendieck-relay.onrender.com;
% PDF page k = archivists' page 40+k, no cover sheet), per the skill
% transcribe-grothendieck. The folder is 84 pages in five batches, done in
% parallel; this is the third.
%
% Pass: Opus 5.5 (claude-opus-5-5), 2026-09-26 — first pass, unchecked against the pages by a human.
%
% Pages skipped: 52, 54 and 58, leaves of an English typescript on invariant
% theory (typed pagination « -1.5- », « -1.6- », « -0.26 bis- »: Definition 1.5
% (Reynolds operator), Theorem 1.1 (universal categorical quotient of an affine
% scheme by a reductive group), Definition 0.10 (principal fibre bundle); it
% cites « Cartier [6] » and « Iwahori [ ] »). This is evidently a draft of
% Mumford's Geometric Invariant Theory, chapters 0-1; it is not his and these
% three leaves carry nothing of his but a pencil box round « §2. » (p. 52) and
% the typed-over page number (p. 58). Skipped by the settled rule, like p. 37
% in batch 2. Pages 46 and 59 are covers (a tan sheet with only his title): no
% \page marker, his words become the section heading, with a \note.
%
% Pages noted only (someone else's typescript, not re-set): 48, 50, 56.
%  - 48 and 50: two leaves of a French typescript, numbered 15 and 13 top
%    right, on Čech cohomology of presheaves of A-modules on a site C
%    (formulas (10)-(14), (21)-(23), Propositions 2.2, 2.3, references
%    « (2.3.14.II) », « [T.F.] Chap. II, §5, Th. 10.2 »). Not an SGA 7 text,
%    not in his voice; author not identified here. One \note each, with the
%    ink interventions in full (hand not established).
%  - 56: a leaf of the same Mumford typescript (« -0.27- », Proposition 0.9
%    and its proof, scanned upside down) carrying pencil comments in English in
%    the margin (« Seems to hold for any equivalence relation ... Is flat
%    enough?? », « Doesn't seem enough, G smooth/k! », « ambiguous »), which
%    are mathematics and presumably his; one \note giving them in full.
%
% Typescript of his (transcribed in full): page 60, the first leaf of a French
% typescript in blue ribbon with his ink corrections, headed by the cover of
% page 59 (« "Th. du cycle invariant" : Cas du H^1 », his hand), which
% continues on page 61 (batch 4). DECISION: transcribed in full as an
% unpublished typescript of his, with the same apparatus as batch 1's
% Griffiths typescript. Evidence: (a) it is in the folder he titled « SGA 7
% (ma part) », under a cover in his hand; (b) his ink corrections are
% mathematical and substantive (« on a d^t = p^h d, où p^h = pgcd(mu_i) »
% replacing a struck typed sentence, the margin « donc h = 0 i.e. d = d_t si k
% est parfait », « projectif », « dfn »), i.e. an author revising, not a reader;
% (c) it speaks of « un résultat de M. RAYNAUD » in the third person, so it is
% not Raynaud's; (d) its numbering — formulas (1)-(10), an unnumbered
% « Théorème », « Corollaire 9 » on p. 61 — does not follow the decimal
% numbering of the published SGA 7 exposés (IX is cited in the literature as
% IX 1.0.3, 2.4, 3.5), and the text is not in batch 1's or batch 2's
% identifications of typed states of VIII and IX. CAVEAT: this pass could not
% compare it line by line with the published Exposé IX (« Modèles de Néron et
% monodromie »), whose §§ on the Néron model of a Jacobian (Raynaud's
% theorem, d and d^t) are close in subject; if it turns out to be a state of a
% published text, it should be reduced to one \note, as in batch 2.
% transcripts/36/batch-04.fr.tex did not exist when this batch was written;
% batches 1, 2 and 5 were read (1 and 2 for their typescript decisions: an
% unpublished typescript of his is given in full, published SGA 7 leaves get
% one note per page; 5 has no typescript).
%
% Pages transcribed from his hand: 41-45, 47, 49, 51, 53, 55, 57 (blue ink on
% bluish or tan paper for 41-45, blue on white for 47-57), and 60 (typescript,
% above). Page 41 was scanned upside down and is cancelled by two long
% diagonals; the lower half of p. 47 is cancelled by diagonals too. Both are
% transcribed, the cancellation recorded once.
%
% His pagination: « 17 » (reading doubtful) top right of p. 41, and nothing
% else in his hand on these leaves. Numbered runs: Définition 1 (p. 43,
% abandoned); a)-b) of p. 47 and a)-b) of p. 49; Corollaire 1 (p. 47),
% Corollaire 2 (p. 49), with a boxed « 2 NB »; a)-b)-c) on pp. 51 and 53
% (a) and b) on p. 51, c) on p. 53, which also re-uses a), b), c) for the
% three conditions on torsors); (*), (**), (***) on p. 47; the typed
% (1)-(5) and a), b) on p. 60. The other typescripts' paginations are given
% in their notes (13, 15; -0.27-).
%
% What the batch is about. (a) p. 41: a birational proper morphism f : X -> Y
% of integral preschemes, U the open set where f^{-1} is defined (complement
% of codim >= 2 when Y is regular in codim 1), and the transport of a divisor
% class xi in Pic(X) to g^{-1}(xi) in Pic(U). (b) pp. 42-43, continuing batch
% 2's « Picard et dualité »: H^1(X, G) = Hom(D(G), Pic_X), H^i(X, D(G)) and
% Ext^i(G, Pic_X), the adjunction Ext^i_X(f^*G, H) = Ext^i_Y(G, f_!H) with
% H = G_m and Pic^•_{X/Y}, the low-degree terms of the resulting spectral
% sequence (i = 0, 1, 2); an abandoned « Foncteurs adjoints, Définition 1 »;
% H^1(X, Hom(gamma, G)) and the exact sequence e -> gamma -> Hom_S(X, gamma)
% -> Hom(X, a; gamma) -> 0. (c) pp. 44-45: the fundamental group of an abelian
% variety A in char. p — p-part via Artin-Schreier-Witt, pi_1^p(A) = T_p(A)
% and the duality T_p(A) x H^1(A, W)^(p) -> Z_p; the prime-to-p part,
% pi_1^l(A) = T_l(A), and via Kummer the Weil pairing T_l(A) x T_l(A*) ->
% T_l(U_l), perfect, and its finite levels. (d) « Biextensions de faisceaux de
% groupes » (cover p. 46), pp. 47-57: Cor. 1-2, describing biextensions of
% P x Q by G through a biextension E' of the subobjects and a pairing phi on
% the T_P, as an equivalence of categories, under divisibility hypotheses;
% « Biextensions sans hyp. de rigidité » (p. 51): Ext^1(A, G) = H^0(S,
% Ext^1(A,G)) when Hom_gr(A,G) = 0, H^1(A, G_A) = H^1(S,G) x H^0(S, R^1f_*G_A)
% when Hom_pct(A,G) = 0, the exact sequence 0 -> Ext^1(A,G) -> R^1f_*(G_A) ->
% Tors birig sym(A,A;G) and NS(A); p. 53, Biext(A,B;G) = Hom(B, Ext^1(A,G))
% and its embedding in rigidified torsors on A x_S B; pp. 55, 57: the inverse
% map Delta^* by the diagonal, Delta^* j(L) = (2 id)^*(L) L^{-2}, the lemma
% (n id)^* L = n^2 L and the quadraticity of Hom(A,B) -> Hom(NS(B), NS(A)),
% whence the canonical extension of NS_{A/S} splits canonically (factor 2).
% (e) p. 60: the invariant-cycle problem for H^1 (after Raynaud): S strictly
% local trait, X projective regular over S with O_S = f_*O_X, X_0 = sum d_i D_i,
% d = pgcd(d_i), Gamma, w : Gamma -> NS(X_0), d^t; a) d^t = pgcd(mu_i d_i),
% d^t = p^h d; b) Ker w is in the torsion of Gamma, cyclic of order dividing d.
%
% Notation fixed once: \underline{\operatorname{Pic}}, \gamma and
% \underline{\mathcal{H}} as in batch 2; \mathcal{E}xt and \mathcal{H} for his
% script Ext and H on pp. 42-43; \mathbb{W} for his double-struck W (Witt
% vectors) on p. 44; \mathbb{P} for his double-struck P (a set of primes) on
% pp. 47-49; {}_nP for the n-torsion, written with the n in front; G_A, e_A,
% e_{A \times_S B} as written; \underline{NS}; typed underlined O and Z as
% \underline{O}, \underline{\mathbf{Z}}.
%
% Apparatus: 98 \ill{}, 70 \uncertain{}. Hardest pages: 47 and 49 (fast notes, the
% cancelled half of 47, the boxed NB of 49), 53 and 55 (connective prose), 41.
% Readings to check first: the bracket of p. 41 (« i.e. aussi ... »); the
% annotation « pt singulier » and the i = 2 line of p. 42; the right-hand
% phrases of p. 44; « Si on admet que ... pas de torsion » on p. 45; the sign
% in the quadratic identity of p. 55; « splitte canoniquement » on p. 57.

\documentclass[11pt,a4paper]{article}
% Le préambule commun à toutes les transcriptions du fonds Grothendieck.
%
% Il tient en une idée : **l'incertitude se compose, elle ne se résout pas.**
% Une transcription qui lisse un mot douteux a détruit la seule chose pour
% laquelle on la faisait — un lecteur doit pouvoir savoir, à chaque endroit,
% ce qui était sur le feuillet et ce qui a été ajouté. Les six macros qui
% suivent sont donc l'appareil critique, et non de la décoration.
%
% Les mêmes commandes sont reconnues par `scripts/render.mjs`, qui produit la
% vue de lecture du site. Une macro ajoutée ici doit l'être aussi là-bas,
% faute de quoi le rendu échouera bruyamment — ce qui est voulu : un rendu
% silencieusement faux est pire qu'un rendu absent.

\ProvidesPackage{grothendieck}[2026/08/08 Transcriptions du fonds Grothendieck]

\usepackage{amsmath,amssymb,amsthm}
\usepackage{mathrsfs}
\usepackage{stmaryrd}
% Les diagrammes commutatifs sont partout dans ce fonds : c'est la langue
% même de la géométrie algébrique de Grothendieck, et une transcription qui
% les rendrait en prose ne transcrirait rien.
\usepackage{tikz-cd}
% Français par défaut — c'est la langue des feuillets — mais l'anglais est
% chargé aussi : la traduction et le résumé sont des documents anglais, et la
% typographie française (espaces avant les deux-points) y serait une faute.
% Ces documents basculent par \selectlanguage{english} en tête de corps.
\usepackage[english,french]{babel}
\usepackage{fontspec}
\usepackage{geometry}
\geometry{a4paper,margin=25mm,marginparwidth=28mm}
\usepackage{marginnote}
\usepackage{xcolor}
% ulem, not soul. `soul`'s \st and \ul are built on a character-by-character
% scan that XeLaTeX's font machinery breaks: the document fails with an
% undefined control sequence rather than degrading. ulem does the same two jobs
% and is Unicode-engine safe. `normalem` stops it hijacking \emph, which the
% transcriptions use for Grothendieck's own emphasis.
\usepackage[normalem]{ulem}
\usepackage{hyperref}
\hypersetup{colorlinks=true,linkcolor=black,urlcolor=black,pdfborder={0 0 0}}

% --- Métadonnées du lot -----------------------------------------------------
% Reprises telles quelles par le script de rendu, qui les affiche en tête de
% la vue de lecture. Elles fixent ce que le fichier prétend couvrir : sans
% elles, une transcription flotte sans point d'attache dans un fonds de seize
% mille feuillets.

\newcommand{\folder}[1]{\gdef\grothFolder{#1}}
\newcommand{\batch}[1]{\gdef\grothBatch{#1}}
\newcommand{\leaves}[2]{\gdef\grothFirst{#1}\gdef\grothLast{#2}}
\newcommand{\foldertitle}[1]{\gdef\grothFoldertitle{#1}}
\gdef\grothFolder{?}\gdef\grothBatch{?}\gdef\grothFirst{?}\gdef\grothLast{?}\gdef\grothFoldertitle{}

% --- Le résumé --------------------------------------------------------------
%
% Il ouvre la lecture modernisée, sous le titre « Résumé », et il est composé
% en petit. Ce n'est pas un ornement : c'est le seul passage du document qui
% ne soit pas des mathématiques, et un lecteur doit pouvoir le reconnaître
% avant de l'avoir lu — puis le sauter s'il connaît déjà le sujet, ou s'y
% arrêter s'il ne le connaît pas. Le filet qui le suit marque la reprise du
% corps.

\newenvironment{resume}
  {\small\setlength{\parskip}{0.45em}}
  {\par\medskip\hrule\medskip}

% --- Mots-clés ---------------------------------------------------------------
%
% En anglais, en clôture du résumé de la lecture modernisée : le vocabulaire
% moderne sous lequel chercher ce que les feuillets construisent. Le manifeste
% du site (`npm run manifest`) les extrait pour étiqueter la cote entière —
% cette ligne est la source unique des tags, il n'y a pas de fichier à côté.
\newcommand{\keywords}[1]{\par\smallskip\noindent
  {\footnotesize\textsc{Keywords} --- \textit{#1}}\par}

% --- L'appareil critique ----------------------------------------------------

% Le numéro de feuillet, en marge. C'est l'ancre qui permet de régler un
% désaccord sur une formule en regardant *un* feuillet plutôt qu'une chemise
% de six cents. Le site s'en sert aussi pour tourner les pages du fac-similé
% quand on fait défiler la transcription.
\newcommand{\leaf}[1]{%
  \marginnote{\footnotesize\color{gray!70}\textsf{#1}}%
  \label{leaf:#1}\ignorespaces}

% Illisible : on ne devine pas. Un mot inventé qui se lit comme les autres est
% le pire résultat possible.
\newcommand{\ill}{\textcolor{red!60!black}{[\,\ldots\,]}}

% Lecture douteuse : le mot est donné, et signalé comme incertain.
\newcommand{\uncertain}[1]{\uline{#1}}

% Ajout de l'éditeur — une lettre manquante, un mot sous-entendu.
\newcommand{\add}[1]{\textcolor{blue!50!black}{[#1]}}

% Biffé par Grothendieck lui-même. Ce qu'il a rayé fait partie du document :
% une rature dit souvent où la pensée a changé de direction.
\newcommand{\struck}[1]{\sout{#1}}

% Note du transcripteur, sur l'état matériel du feuillet.
\newcommand{\note}[1]{\par\smallskip{\small\color{gray!60!black}\textit{[#1]}}\par\smallskip}

% Note marginale de Grothendieck, distincte du corps du texte.
\newcommand{\marginal}[1]{\par\smallskip\noindent\hspace{1em}%
  {\small\color{gray!60!black}#1}\par\smallskip}

% --- Titre ------------------------------------------------------------------

\AtBeginDocument{%
  \begin{center}
    {\large\bfseries Fonds Alexandre Grothendieck --- cote n\textsuperscript{o}~\grothFolder}\\[2pt]
    {\itshape\grothFoldertitle}\\[4pt]
    {\small feuillets \grothFirst--\grothLast\ (lot \grothBatch)}\\[2pt]
    {\footnotesize\color{gray!60!black}Transcription établie d'après le fac-similé numérisé
     par l'Université de Montpellier.\\
     Les lectures incertaines sont soulignées, les illisibles notées [\,\ldots\,],
     les ajouts entre crochets.}
  \end{center}
  \vspace{1em}}

\endinput


\folder{36}
\batch{3}
\pages{41}{60}
\dating{[vers 1967-1973]}
\watermark{Édition de démonstration}
\foldertitle{SGA 7 (ma part) : notes manuscrites (s.d.), tapuscrit annoté (s.d.)}

\begin{document}

\section*{Picard et dualité (suite)}

\note{le titre de section est repris du feuillet de garde de la page 34
(lot précédent) ; rien n'indique sur les feuillets qui suivent qu'ils lui
appartiennent tous}

\page{41}

\note{toute la page est barrée de deux longues diagonales ; elle se lit sous
elles. En haut à droite, un numéro \uncertain{17}, de sa main}

\begin{tikzcd}
f^{-1}(U) \arrow[r, hook, "i"] \arrow[d, "f"'] & X \arrow[d, "f"] \\
U \arrow[r, "j"'] \arrow[u, bend right=30, "g"'] & Y
\end{tikzcd}

\[
gf = i, \qquad fg = j
\]
\note{le schéma et les deux relations sont dans la marge de gauche ; la flèche
$g$ monte de $U$ vers $f^{-1}(U)$ en diagonale}

$f$ morphisme birationnel \add{propre} de préschémas intègres $X$, $Y$.
\note{« propre » est ajouté au-dessus de la ligne}

$U$ ouvert de $Y$ sur lequel $f^{-1}$ est défini [si $Y$ est régulier en
codim.~1, alors $U$ est le complémentaire d'une \uncertain{partie} de codim.\
$\geqslant 2$], \struck{i.e.\ \uncertain{aussi les} \ill{}} i.e.\ \uncertain{aussi
les} \ill{} \uncertain{ont un} voisinage $V$ tel que $f$ induise un isom.\
$f^{-1}(V) \to V$.

Soit $\xi \in \operatorname{Pic}(X)$ défini par un \struck{\uncertain{cycle}}
\emph{\uncertain{divisoriel}} $D$ \uncertain{(Cartier)}, considérons
\struck{\ill{}} \struck{le morphisme} $g^{-1}(\xi) \in \operatorname{Pic}(U)$, il
\uncertain{correspond} à un \struck{\ill{}} cycle divisoriel
\note{le mot rayé au-dessus de « divisoriel » est relié par un trait à
l'endroit de l'insertion ; la page s'arrête sur « cycle divisoriel »}

\page{42}

\[
\mathcal{H}^1(X; A) \qquad \mathcal{E}xt^1(A, \underline{\operatorname{Pic}})
\]
\note{sous $\mathcal{H}^1(X;A)$, un trait descend vers l'annotation
« \uncertain{pt singulier} » ; à gauche, deux croquis de courbes : une boucle
fermée marquée d'un point, et une courbe qui se recoupe en un point marqué}

\[
\boxed{\mathcal{H}^1(X, G) \simeq \operatorname{Hom}(D(G),
\underline{\operatorname{Pic}}_X)}
\]
\[
\mathcal{H}^i(X, D(G)) \simeq \operatorname{Ext}^i(G,
\underline{\operatorname{Pic}}_X)
\]
\[
\mathcal{E}xt^i_X(f^*(G), \mathbb{G}_m) \simeq \mathcal{E}xt^i_Y(G,
f_{\cdot}(\mathbb{G}_m))
\]
\note{l'indice de $f$ au second membre est un simple point}

\[
\boxed{\underline{\mathcal{E}xt}^i_X(f^*(G), H^{\cdot}) \simeq
\underline{\mathcal{E}xt}^i_Y(G, f_!(H^{\cdot}))}
\qquad \text{avec } f\colon X \to Y
\]
\note{la flèche verticale $f\colon X \to Y$ est dessinée dans la marge, à gauche
de l'encadré}

$H = \mathbb{G}_m$
\[
\underline{\mathcal{E}xt}^i_X(f^*(G), \mathbb{G}_m) \simeq
\underline{\mathcal{E}xt}^i_Y(G, \underline{\operatorname{Pic}}^{\bullet}_{X/Y})
\]
\note{le point en exposant de $\underline{\operatorname{Pic}}$ est cerclé et
annoté « les \uncertain{supérieurs} » ; sous le premier membre, une double
flèche montante part de « \struck{\ill{}}$(X,$ », et sous le signe
$\simeq$ un double trait vertical descend vers la ligne suivante}

[$G$ \emph{\uncertain{fini}}, alors
\quad $H^i(X, D(G)) \Leftarrow$ \struck{\ill{}} $\mathcal{E}xt^i(G,
\underline{\operatorname{Pic}}^{(1)}_{X/Y})$
\[
\begin{array}{ll}
i = 0 & H^0(X, D(G)) = D(G) \\
i = 1 & H^1(X, D(G)) = \operatorname{Hom}(G, \underline{\operatorname{Pic}}^{(1)}) \\
i = 2 & H^2(X, D(G)) \neq \mathcal{E}xt^1(G, \underline{\operatorname{Pic}}^{(1)}_{X/Y})
+ \operatorname{Hom}(G, \underline{\operatorname{Pic}}^{(2)}_{X})
\end{array}
\]
\note{la ligne entre crochets est encadrée ; devant $D(G)$ (cas $i=0$) un
« Hom » est biffé, devant $\operatorname{Hom}$ (cas $i=1$) un « Ext » ; sous le
$\neq$ du cas $i=2$, un point d'exclamation ; le dernier terme, au bord droit,
est lu en partie. L'exposant $(1)$ est lu sous réserve}

\ill{} \ill{} \ill{} \quad $H^{i-1}(X, G') \Leftarrow \mathcal{E}xt^{i}(G,
\underline{\operatorname{Pic}}^{?}_{X/Y})$
\note{la dernière ligne, au pied de la page, est lue en partie ; les exposants
du second membre sont illisibles (un point d'interrogation, peut-être le sien,
tient lieu de l'exposant de $\underline{\operatorname{Pic}}$). Au-dessous :
« $i=1$ » et « $H^0(X, G') =$ », sans suite}

\page{43}

\emph{Foncteurs adjoints}

\emph{Définition 1} Soient $X, Y \in (\mathrm{Cat})$, $f\colon X \to Y$ et
$g\colon Y \to X$, on dit que $(f, g)$ est un couple de foncteurs adjoints, ou
que $g$ est adjoint de $f$, ou que $f$ est coadjoint de $g$, s'il existe un
\struck{\ill{}}
\note{la définition s'arrête là ; un long trait oblique la sépare du reste de
la page}

\[
\mathcal{H}^1(X, \underline{\operatorname{Hom}}(\gamma, G)) \simeq
\mathcal{H}^1(X, \underline{\operatorname{Ext}}^1(\gamma, G))
\]
\note{écrit sur deux lignes, le signe $\simeq$ en fin de première ligne ; en
dessous, « $\operatorname{Ext}^1(\gamma,$ » commencé puis abandonné sous deux
traits. Même $\gamma$ bouclé qu'à la page 40}

\[
X \to \underline{\operatorname{Hom}}_{S\text{-gr}}(\gamma, G), \qquad
X \times \gamma \to G, \qquad
\gamma \to \underline{\operatorname{Hom}}_S(X, G)
\]
\note{écrit en colonne à droite ; sous le troisième, un croquis relie $S$ et
$G$ par un trait marqué $\mathrm{pr}_2$, puis « $\operatorname{Hom}(\gamma,$ »
inachevé}

\[
\prod^{\circ}_{X/S} \gamma_X = \underline{\operatorname{Hom}}_S(X, \gamma)
\doteq \gamma
\]
\[
e \to \gamma \to \underline{\operatorname{Hom}}_S(X, \gamma) \to
\operatorname{Hom}(X, a; \gamma) \to 0
\]
\note{le premier membre est écrit sur une ligne et le second, précédé d'un
double trait vertical, sous lui ; le petit cercle est au-dessus du
$\prod$}

\page{44}

D'après la théorie d'Artin-Schreier-Witt, on a
\[
H^1(A, \mathbb{Z}/p) \simeq H^1(A, \underline{O}_A)^{(p)}
\qquad \text{éléments de } H^1(A, \underline{O}_A) \text{ tels que } Fx = x
\]
\[
H^1(A, \mathbb{Z}/p^n) \simeq H^1(A, \underline{O}_A(W_n))^{(p)}
\qquad H^1(A, \underline{O}_A(W_n)) \text{ tels que } Fx = x
\]
\[
\boxed{\operatorname{Hom}(\pi_1^p(A), \mathbb{Z}_p) = H^1(A, \mathbb{Z}_p) \simeq
H^1(A, \mathbb{W})^{(p)}}
\]
\ill{} \ill{} \uncertain{dualité entre}
\note{à droite de l'encadré, quelques mots rapides lus en partie ; dans
l'encadré, un signe raturé devant $\mathbb{Z}_p$, et un autre en indice
de la parenthèse finale}

\struck{accouplement} \uncertain{Donc explicitons}

\struck{$\pi_1^p(A) \times H^1(A, \mathbb{W})_s \to \mathbb{Z}_p$}

qui est une dualité \struck{\ill{}} de $\mathbb{Z}_p$-modules libres de type fini
\[
\boxed{\operatorname{Hom}_{\mathrm{cont}}(\pi_1^p(A), \Lambda) \simeq H^1(A,
\mathbb{W})_s}
\]
\note{la formule biffée et la suivante sont encadrées ensemble et reliées par
un grand arc ; le $\Lambda$ est lu sous réserve}

D'autre part, utilisant \ill{} le fait que \uncertain{tout} revêtement de $A$ est
une variété abélienne, on trouve \uncertain{aussi}
\[
\boxed{\pi_1^p(A) \simeq T_p(A)}
\]
\uncertain{d'où} \ill{} \uncertain{accouplement}
\[
\boxed{T_p(A) \times H^1(A, \mathbb{W})^{(p)} \longrightarrow \mathbb{Z}_p}
\]
qui est une dualité de $\mathbb{Z}_p$-modules libres de type fini.

\struck{qu'on peut aussi écrire} \uncertain{dans} \ill{} \uncertain{dualité}
\[
\boxed{\operatorname{Hom}_{\mathrm{cont}}(T_p(A), W_\infty(k)) \simeq H^1(A,
\mathbb{W})_s}
\]
\note{devant $W_\infty(k)$, une lettre raturée}

\page{45}

$A$ variété abélienne en car.\ $p$. On a, pour la partie de $\pi_1(A)$ première à
$p$
\[
\pi_1^{\neq p}(A) \simeq \varprojlim_{\substack{n \\ (n,p)=1}} A^{(n)}
\]
où $A^{(n)}$ est le noyau de $A \xrightarrow{\ n\ } A$, et où $A^{(mn)} \to A^{(n)}$
est la multiplication par $m$.
\[
\boxed{\pi_1^{\ell}(A) \simeq \varprojlim_{k} A^{(\ell^k)} \simeq T_\ell(A) =
\operatorname{Hom}(\mathbb{Q}_\ell/\mathbb{Z}_\ell, A)}
\qquad \text{module de Tate}
\]
$\ell$ étant un nb premier $\neq p$. Ce résultat provient du fait qu'un
\emph{revêtement non ramifié d'une variété abélienne est une variété abélienne}.

D'après Kummer, on a d'autre part
\struck{\ill{}}
\[
\operatorname{Hom}(\pi_1^{\ell}(A), U_\ell) \simeq
\text{sous-groupe de } \underline{\operatorname{Pic}}(A) = H^1(A,
\underline{O}_A^*)
\]
des éléments d'ordre une puissance de $\ell$.

Si on admet que $A$ \ill{} \emph{\uncertain{pas de torsion}}, on obtient donc
\[
\operatorname{Hom}(\pi_1^{\ell}(A), U_\ell) \simeq A^{*(\ell^\infty)}, \quad
\text{i.e.}
\]
\[
\boxed{T_\ell(A) \times T_\ell(A^*) \longrightarrow T_\ell(U_\ell)}
\]
accouplement à valeurs dans le $\mathbb{Z}_\ell$-module libre de rang $1$
$T_\ell(U_\ell)$, et accouplement qui est une dualité.
\[
T_\ell(A)/\ell^k T_\ell(A) \times T_\ell(A^*)/\ell^k T_\ell(A^*) \longrightarrow
T_\ell(U_\ell)/\ell^k T_\ell(U_\ell)
\]
\[
A^{(\ell^k)} \times A^{*(\ell^k)} \longrightarrow U^{(\ell^k)}
\]

\section*{Biextensions de faisceaux de groupes}

\note{titre inscrit au crayon en haut à droite du feuillet de garde
(page~46), qui ne porte rien d'autre ; le feuillet ne reçoit pas de numéro de
page}

\page{47}

\[
0 \to P' \to P \to P'' \to 0 \qquad Q
\]
\struck{\ill{}} a) ${}_nQ = Q$ \quad b) ${}_nP \to P''$ \uncertain{épimorphisme}

Toute biextension \add{$E$} de $P \times Q$ \struck{définit} par $G$ définit

a) Une biextension \add{$E'$} de $P' \times Q$ par $G$, d'où un accouplement
\[
{}_nP' \times {}_nQ \xrightarrow{\ E'(n)\ } G
\]

b) Une biextension de ${}_nP \times Q$ par $G$, i.e.\ un accouplement
\[
{}_nP \times {}_nQ \xrightarrow{\ \varphi\ } G
\]
\note{$E$ et $E'$ sont ajoutés au-dessus de la ligne. À droite de a) et b),
un bloc accolé et rayé de hachures obliques, lu en partie :
\struck{Le \ill{} $0 \to P' \to P \to P'' \to 0$ \ill{} \ill{}
$\operatorname{Hom}(P\otimes Q, G) \to \operatorname{Hom}(P'\otimes Q, G)$
\ill{} \ill{} ; $\underline{\operatorname{Hom}}(P'\otimes Q, G) = 0$}}

\struck{Avec les} Ces deux $(E', \varphi)$ satisfont à la condition que
\[
\varphi|({}_nP' \times {}_nQ) = E'(n)
\]

Un hom.\ de biextensions $E_1 \to E_2$ définit un hom.\ $(E'_1, \varphi_1) \to
(E'_2, \varphi_2)$, i.e.\ \add{\uncertain{(déf)}} un hom.\ de biextensions
$u\colon E'_1 \to E'_2$, tel que $\varphi$ \ill{}.

Je dis que le \emph{foncteur} $E \mapsto (E', \varphi)$ est une \emph{équivalence
de catégories}.

\struck{a) \ill{}} \struck{a) $\operatorname{Hom}(P\otimes Q, G) \simeq
\operatorname{Hom}(P'\otimes Q, G)$} \quad \struck{b) $E_1 \simeq E_2$
\ill{} $(E'_1, \varphi_1) \simeq (E'_2, \varphi_2)$}

\emph{Pl.\ fidèle} signifie aussi que $E \mapsto E|(P', Q)$ est pl.\ fid., i.e.\
\ill{} \ill{}, que $\underline{\operatorname{Hom}}(P''\otimes Q, G) = 0$,
\uncertain{(ou)} $\operatorname{Hom}(P\otimes Q, G) \to \operatorname{Hom}(P'\otimes
Q, G)$ surj.\ \struck{(**) \ill{} \ill{} tel que $(E', \varphi) \simeq$ \ill{}
\ill{}, si $E$ est tel que $E' = 1$, $\varphi = 0$, alors $E = 1$} La
première condition résulte de a) et b) qui impliquent $P''\otimes Q = 0$, la
deuxième résulte de c) [dans la deuxième alternative de c) les deux
\uncertain{côtés} de (**) sont nuls !].

\struck{Exemple}. Pour (***), on note que si $E' = 0$, \ill{} \ill{} la suite
exacte, $E$ est déterminé par un élément de $\operatorname{Ext}^1(P''\otimes Q,
G)$, ou aussi équivalemment par un élément de $\operatorname{Hom}(P''\otimes_{?} Q,
G)$, on \uncertain{donne} que c'est ce que définit $\varphi$ \ill{} \ill{}.
\note{toute la partie de la page qui va de « Pl.\ fidèle » à « ce que définit
$\varphi$ » est barrée de longues diagonales ; elle se lit sous elles. Le
c) invoqué n'est pas sur la page}

\emph{Corollaire 1} Supposons, pour un $\mathbb{P}$ \ill{} \ill{}, que
$(\varinjlim_{n} {}_nP) \to P''$ surjectif, et que $p \in \mathbb{P} \Rightarrow
pQ = Q$. Alors on a une équivalence $E \mapsto (E', \varphi)$, où $E'$ est une
\note{sous la limite, « \uncertain{$n$ décrivant} $\mathbb{P}$ » ; le texte
continue page 49}

\page{48}

\note{page non retranscrite : feuillet d'un tapuscrit en français, numéroté 15
en haut à droite, qui n'est pas de lui et n'appartient pas à SGA 7 (préfaisceaux
et faisceaux de $A$-modules sur un site $C$ : le foncteur faisceau associé
$a$, adjoint à gauche de $\mathcal{H}^0$, formules (21)--(23), Propositions 2.2
et 2.3 liant la cohomologie de Čech à celle des faisceaux, renvois
« (2.3.14.II) » et « [T.F.] Chap.~II, §~5, Th.~10.2 »). Interventions à
l'encre, d'une main non établie ici : « \uncertain{adjoint} » entouré dans la
marge gauche, en regard d'un passage tapé entre crochets ; en marge, la formule
$F(aS) \xrightarrow{\sim} \mathcal{H}^0(F)(S)$ « i.e. : », appelée devant
la formule (22), dont une première version tapée est biffée ; « Laisser ici un
interligne », entouré, avec une flèche, avant « Nous allons maintenant
énoncer » ; des crochets devant les deux propositions ; « fonctorielle »
corrigé.}

\page{49}

biextension \uncertain{restreinte} de $P, Q$ par $G$, et où $(E', \varphi)$ est le
couple d'une biext.\ $E'$ de $P', Q$ par $G$ et d'un hom.
\[
\varphi\colon T_{\mathbb{P}}(P) \times T_{\mathbb{P}}(Q) \longrightarrow
T_{\mathbb{P}}(G)
\]
qui \uncertain{étende} l'hom.
\[
\varphi_{\mathbb{P}, E'}\colon T_{\mathbb{P}}(P') \times T_{\mathbb{P}}(Q)
\longrightarrow T_{\mathbb{P}}(G)
\]
\note{devant $\varphi_{\mathbb{P},E'}$, une lettre biffée}

\marginal{2 NB On considère les $T_{\mathbb{P}}(P)$ comme \uncertain{syst.}\
\uncertain{projectifs} \struck{\uncertain{stricts}} \ill{} \ill{} \ill{}
\uncertain{syst.\ projectif} \ill{} \ill{} \uncertain{limite}) ; cette description
\ill{} \ill{} si $\forall p \in \mathbb{P}$, \ill{} ${}_pP' = P'$, \ill{} ; il faut
garder le syst.\ \ill{} ${}_nP \times {}_nQ \to {}_nG$ \struck{\ill{}} avec la
compatibilité \ill{} \ill{} des restrictions \ill{} ${}_{n}P' \times {}_{n}Q$
\ldots}
\note{le « 2 NB » est encadré dans la marge droite, à hauteur de la formule
précédente et du corollaire 2 ; écrit serré avec plusieurs surcharges, il
n'est lu qu'en partie}

\emph{Corollaire 2} Soient deux suites exactes
\[
0 \to P' \to P \to P'' \to 0, \qquad 0 \to Q' \to Q \to Q'' \to 0
\]
et supposons que, pour un $\mathbb{P}$ de nbs premiers, on ait :

a) $p \in \mathbb{P} \Rightarrow p P = P$, $p Q' = Q'$ \struck{\ill{}} (ou le
symétrique) \struck{(\uncertain{ou} $pP = P$ \ill{} $pQ = Q$)}

b) \uncertain{Les} $\varinjlim_{n \to \mathbb{P}} {}_nP \to P''$ et
$\varinjlim_{n \to \mathbb{P}} {}_nQ \to Q''$ épim.

\marginal{NB \uncertain{remarquer} que, \ill{} \ill{} tel que \ill{} \ill{} que
$P', Q', P'', Q''$ soient \ill{} $\mathbb{P}$ \ill{} dans a), \ill{} \ill{} de b),
\ill{} ${}_nP \times {}_nQ \to {}_nG$ \ldots}
\note{second NB encadré, à droite de a)--b), lu en partie}

Alors le foncteur
\[
E \longmapsto (E', \varphi)
\]
\uncertain{envoyant} : \uncertain{les} biext.\ de $P, Q$ par $G$ \uncertain{sur}
\ill{} \ill{} restriction $E'$ à $P', Q'$ et la forme $\varphi_{E,\mathbb{P}}$
\[
\varphi\colon T_{\mathbb{P}}(P) \times T_{\mathbb{P}}(Q) \longrightarrow
T_{\mathbb{P}}(G)
\]
prolongeant
\[
\varphi_{E';\mathbb{P}}\colon T_{\mathbb{P}}(P') \times T_{\mathbb{P}}(Q')
\longrightarrow T_{\mathbb{P}}(G)
\]
est une équivalence de catégories.

\emph{Dém.} On applique d'abord le \struck{\ill{}} corollaire 1 à
$0 \to P' \to P \to P'' \to 0$ et $Q'$, en utilisant $pQ' = Q'$,
$\varinjlim_{n} {}_nP \to P''$ épim. Cela donne une description des
biextensions de $P, Q'$. On applique ensuite le cor.~1 à $P$ et $0 \to Q' \to Q
\to Q'' \to 0$, en utilisant $pP = P$ \struck{\ill{}} et $\varinjlim_{n\to
\mathbb{P}} {}_nQ \to Q''$ épim.
\note{sous « équivalence de catégories », une petite ligne,
« \uncertain{supposons} \ill{} », et au-dessus de « corollaire 1 », un mot
biffé}

\page{50}

\note{page non retranscrite : feuillet du même tapuscrit, numéroté 13 en
haut à droite (le foncteur $\mathcal{H}^0\colon C^{\sim}_A \to C^{\wedge}_A$ et ses
dérivés $\mathcal{H}^q$, formules (10)--(14), préfaisceaux de cohomologie de
Čech $\check{\mathcal{H}}^q(G)$). Interventions à l'encre : au premier membre
de la formule (12), la version tapée est noircie et remplacée par
$\mathcal{H}^q(F)(X)$ ; de même en (14) par $\check{\mathcal{H}}^q(G)(X)$ ; tout le
passage entre deux traits qui précède (14) est haché de traits obliques ;
« Ainsi qu' » biffé et remplacé par « Comme » ; « les » ajouté dans la marge
avant « préfaisceaux d'$A$-modules »}

\page{51}

\emph{Biextensions sans \uncertain{hyp.}\ de rigidité}

$A$, $G$ faisceaux abéliens sur un topos $\mathcal{S}$, $f\colon A \to S$ \ill{}
\note{le titre est écrit au-dessus de la première ligne et souligné ; les mots
après $f\colon A \to S$ sont en indice, illisibles}

a) Supposons $\underline{\operatorname{Hom}}_{\mathrm{gr}}(A, G) \simeq 0$. Alors
\[
\operatorname{Ext}^1(A, G) \xrightarrow{\ \sim\ } H^0(S,
\underline{\operatorname{Ext}}^1(A, G))
\]
donc la connaissance des \emph{faisceaux} $\underline{\operatorname{Ext}}^1(A, G)$
implique celle de $\operatorname{Ext}^1(A, G)$ et \uncertain{celle} des
$\operatorname{Ext}^1(A_{?}, G_{?})$.
\note{les indices de $A$ et $G$ dans ce dernier terme sont illisibles}

b) Supposons $\underline{\operatorname{Hom}}_{\mathrm{pct}}(A, G) = 0$,
[i.e.\ $f_*(G_A) \xleftarrow{\ \sim\ } G$]. Alors
\[
H^1(A, G_A) \simeq H^1(S, G) \times H^0(S, R^1f_*(G_A))
\]
\note{sous « pct », en petit : « Hom (\uncertain{pas nécessairement}
multiplicatifs) qui respectent la section unité ». Le coefficient du premier
membre est écrit plutôt $A_G$ ; on le rend par $G_A$, comme au
second membre}

\[
\begin{array}{c}
H^0(S, R^1f_*(G_A)) \\
\wr \\
\text{Classes d'isom.\ \uncertain{prises} des } G_A\text{-torseurs rigidifiés
\uncertain{rel.}\ à } e_A \\
\wr \\
\text{Classes de } G_A\text{-torseurs rigidifiables par } e_A.
\end{array}
\]
\note{dans la colonne de droite, sous le dernier facteur ; devant
« rigidifiés », un mot biffé, et devant « $G_A$-torseurs » de la dernière
ligne, « torseurs » biffé}

De plus, le foncteur \uncertain{évident}, qui va de la catégorie des extensions
de $A$ par $G$ dans la catégorie des $G_A$-torseurs $e_A$-rigidifiés, est un
foncteur pl.\ fid.\ (de catégories discrètes), (\struck{\ill{}} on a une
\emph{injection}
\[
0 \longrightarrow \operatorname{Ext}^1(A, G) \longrightarrow H^1(A, G_A) \quad );
\]
l'image essentielle est formée des \struck{\ill{}} $G_A$-torseurs
\struck{\ill{} tels que \ill{} $\xi$ \ill{}}
\[
\pi^*(L) \simeq \mathrm{pr}_1^*(L)\, \mathrm{pr}_2^*(L) \qquad \text{i.e.\ tels
que}
\]
\[
\pi^*(L)\, \mathrm{pr}_1^*(L)^{-1}\, \mathrm{pr}_2^*(L)^{-1} \simeq 0 .
\]
On trouve \ill{} \ill{} \ill{} \uncertain{conclusion} une suite exacte
\[
A \times A \overset{\mathrm{pr}_1,\ \mathrm{pr}_2,\ \pi}{\longrightarrow} A
\]
\note{trois flèches parallèles, marquées $\mathrm{pr}_1$, $\mathrm{pr}_2$,
$\pi$, dans la marge de gauche}

\begin{tikzcd}[column sep=small, row sep=small, nodes={font=\scriptsize}]
0 \arrow[r] & \underline{\operatorname{Ext}}^1(A, G) \arrow[r, "i"] & R^1f_*(G_A) \arrow[r, "j"] \arrow[dl, "k"'] & \underline{\operatorname{Tors}}\,\underline{\operatorname{birig}}\,\underline{\operatorname{sym}}(A, A; G) \arrow[d, no head, "\Vert" description] \\
0 \arrow[r] & \underline{NS}(A) \arrow[rr, "j'"] & & \underline{\operatorname{Tors}}\,\underline{\operatorname{birig}}\,\underline{\operatorname{sym}}(A, A; G)
\end{tikzcd}

\note{la flèche $k$ descend en oblique de $R^1f_*(G_A)$ vers $\underline{NS}(A)$ ;
entre $\underline{NS}(A)$ et $j'$, un signe noirci. À gauche de la seconde ligne,
« i.e. »}

\page{53}

c) Soit $B$ un autre faisceau abélien. Supposons de même que
$\underline{\operatorname{Hom}}_{\mathrm{gr}}(A, G) = 0$. Alors on trouve que la
\struck{\uncertain{fonction}} \add{catégorie} \struck{\uncertain{naturel}}
$\underline{\operatorname{BIEXT}}(A, B; G)$ est discrète et

\struck{$\operatorname{Biext}(A, B; G) \longleftarrow$}

\[
\operatorname{Biext}(A, B; G) \xrightarrow{\ \sim\ }
\underline{\operatorname{Hom}}_{\mathrm{gr}}(B, \underline{\operatorname{Ext}}^1(A,
G))
\]
\[
\operatorname{Biext}(A, B; G) = H^0(S, \underline{\operatorname{Biext}}(A, B; G))
\simeq \underline{\operatorname{Hom}}_{\mathrm{gr}}(B,
\underline{\operatorname{Ext}}^1(A, G))
\]
\note{au-dessus de $\underline{\operatorname{Ext}}^1(A,G)$ dans la première
formule, « $A'$ »}

Si on suppose \add{\uncertain{de plus}} seulement que \struck{Biext}
$\underline{\operatorname{Hom}}_{\mathrm{pct}}(A, G) = 0$, on trouve
\[
\operatorname{Biext}(A, B; G) \hookrightarrow
\underline{\operatorname{Hom}}_{\mathrm{gr}}(B, R^1f_*(G_A)) \hookrightarrow
H^0(B, R^1f_{B*}(G_{A \times_S B}))
\]
\note{sous $R^1f_*(G_A)$, une variante noircie ; au-dessus de la seconde
inclusion, $\underline{\operatorname{Hom}}_{\mathrm{ens}}(B,
R^1f_*(G_A)) \simeq$ le dernier terme, l'indice « ens » lu sous réserve}
\[
\begin{array}{c}
H^0(B, R^1f_{B*}(G_{A \times_S B})) \\
\Vert \\
\text{classes de } G_{A \times_S B}\text{-torseurs rigidifiés par : }
e_{A \times_S B} \\
\downarrow \\
H^1(A \times_S B, G_{A \times_S B})
\end{array}
\]
\note{dans la colonne de droite ; après « $G_{A\times_S B}$-torseurs », deux
lignes biffées. Une accolade y rattache, à gauche : « $=$ classes de
$G_{A\times_S B}$-torseurs rigidifiables par $e_{A\times_S B}$ »}

\begin{tikzcd}[row sep=small]
A \arrow[r, no head] \arrow[d, "f"'] & A \times_S B \arrow[d, "f_B"] \\
S \arrow[r, no head] & B
\end{tikzcd}

Les torseurs sur $A \times_S B$ \struck{\ill{}} par $G_{A \times_S B}$ qui
\uncertain{définissent} des \ill{} sont ceux qui a) sont $e_{A \times_S B}$
rigidifiables ; et pour lesquels l'application $B \to R^1f_*(G_A)$ qui les
définit a deux vertus : \quad b) c'est additif \quad c) \uncertain{ça} se factorise
par
\marginal{\uncertain{Rappel} : les conditions a) et c) ensemble signifient que
l'hom $A \to R^1q_*(G_B)$ est additif}

\struck{\ill{}} $\underline{\operatorname{Ext}}^1(A, G) \simeq A'$.
\note{le rappel est écrit en interligne sous c), rattaché par une accolade ; le
c) s'arrête sur « par »}

\struck{Supposons qu'on ait également} Notons d'ailleurs que la condition a)
implique que $L$ est \emph{birigidifiable} par.\ à $e_A$, $e_B$, et on trouve un
\uncertain{foncteur} \emph{\uncertain{pl.\ fid.}} de catégories discrètes
\[
\underline{\operatorname{BIEXT}}(A, B; G) \longrightarrow
\underline{\operatorname{TORS\,BIRIG}}_{e_A, e_B}(A, B; G) \longrightarrow
\underline{\operatorname{TORS\,RIG}}_{e_{A \times_S B}}(G_{A \times_S B})
\]
La description \struck{en termes des \ill{} \ill{}} \uncertain{ci-dessus} est
toujours \uncertain{vérifiable}, car a) il est \struck{\ill{}} symétrique en $A$,
$B$, b) dans deux cas importants, le \uncertain{pl.\ fidèle} est une
\uncertain{équivalence} de catégories.

\page{55}

Reprenons la suite exacte de la \uncertain{fin} de \uncertain{5)},
\marginal{où $j'$ signifie que les bitorseurs birigidifiés sur $A$, $A$ sont
des biextensions}
on va définir un hom en sens inverse
\[
R^1f_*(G_A) \xleftarrow{\ \Delta^*\ }
\underline{\operatorname{Tors}}\,[\underline{\operatorname{birig}}\
\underline{\operatorname{sym}}](A, A; G)
\]
\marginal{\uncertain{inutile} pour définir $\Delta^*$}
en utilisant l'application diagonale $\Delta\colon A \to A \times_S A$.
\note{la bulle de la première marginale, en haut à droite, est rattachée à
« la fin de \uncertain{5)} » ; la seconde renvoie par une flèche au crochet
autour de « birig sym ». Pour la diagonale, la page écrit « $A \times_S A
\text{---} A$ », sans pointe lisible}

Calculons le composé \struck{\ill{}} \add{$k \Delta^* j'$} $\Delta^* j'$ avec
$k\colon R^1f_*(G_A) \to \underline{NS}(A)$, on aura d'abord
\[
\Delta^* j(L) \simeq (2\,\mathrm{id}_A)^*(L)\, L^{-2}
\]
d'où \struck{Je dis que} \ill{} $\equiv L^2$ \quad $\lbrace \bmod A' =
\underline{\operatorname{Ext}}^1(A, G) \rbrace$.

Cela résulte du fait suivant

\emph{Lemme} Pour tout entier $n$, $(n\,\mathrm{id}_A)^*(L) \equiv n^2 L$,
\struck{\ill{} \ill{}} plus généralement, si $A$, $B$ \ill{} \uncertain{comme
ci-dessus}, \ill{}
\[
\operatorname{Hom}(A, B) \longrightarrow \operatorname{Hom}(\underline{NS}(B),
\underline{NS}(A))
\]
est une application \emph{quadratique}, i.e.
\[
\varphi(u+v) - \varphi(u-v) = 2(\varphi(u) + \varphi(v)) .
\]
\note{le signe entre les deux premiers termes est lu « $-$ » sous réserve ;
l'identité du parallélogramme porterait « $+$ ». Transcrit tel que lu}

\struck{En effet} \emph{Dém.} En effet, cette application se factorise en
\[
\begin{array}{c}
\operatorname{Hom}(A, B) \longrightarrow
\operatorname{Hom}(\underline{\operatorname{Bitors}}(B, B; G),
\underline{\operatorname{Bitors}}(A, A; G))\ ! \\
\downarrow \ \text{linéaire} \\
\operatorname{Hom}(\underline{NS}(B), \underline{NS}(A))
\end{array}
\]
\ill{} \uncertain{s'expliquant} ! indiquer les hom qui envoient
$\underline{NS}(B)$ dans $\underline{NS}(A)$. \struck{\ill{}} Or il est évident que
la première flèche est quadratique.
\note{« linéaire », à côté de la flèche verticale, est lu sous réserve}

\page{56}

\note{page non retranscrite : feuillet d'un tapuscrit anglais qui n'est pas de
lui (numéroté « -0.27- » ; « Proposition 0.9 » : si $\sigma$ est libre et $(Y,
\phi)$ quotient géométrique de $X$ par $G$ sur un corps, $X$ est un fibré
principal sur $Y$, et sa démonstration), de la même série que les pages 52,
54 et 58 non retenues, qui sont d'une rédaction de la théorie géométrique des
invariants. Annotations au crayon, en anglais, présumées de sa main : dans un
cadre en tête de page, « Seems to hold for any equivalence relation $R
\rightrightarrows X$ with "geometric quotient" $Y$, and that $p_1$ be
\emph{smooth}, $X$, $Y$ be \ill{} --- Is \emph{flat} enough ?? » ($p_1$ écrit
sous la double flèche) ; en marge de l'énoncé, « Doesn't seem enough, $G$
\emph{smooth}/$k$ ! », avec un triangle en regard de « $G$, » et un cercle
autour de « is free » ; en marge de la démonstration, « ambiguous », en regard
de « the ideal defining the orbit $O(x)$ », où « the » et « the orbit
$O(x)$ » sont entourés}

\page{57}

Considérons le composé
\[
\Delta^* j\colon \underline{NS}_{A/S} \longrightarrow \underline{P}_{A/S}
\]
on voit que son composé avec $\underline{P}_{A/S} \xrightarrow{\ k\ }
\underline{NS}_{A/S}$ est $2\,\mathrm{id}_{\underline{NS}_{A/S}}$. Donc \emph{la
\ill{} \ill{}} de \emph{l'extension canonique} de $\underline{NS}_{A/S}$
\uncertain{par} $P_{A/S}$ \emph{\uncertain{splitte}} \emph{canoniquement}.
\note{après $\underline{NS}_{A/S}$ et devant $\underline{P}_{A/S}$ dans la première
formule, deux surcharges noircies}

Considérons d'autre part la restriction de $j'(k \Delta^*)$ à
$\underline{\operatorname{Tors}}\,\underline{\operatorname{birig}}\,\underline{\operatorname{sym}}(A,
A; G)$. Je dis que c'est la multiplication par $2$, i.e.
\note{cette phrase est rattachée par une accolade à la ligne précédente ; au
pied, isolé, « $k\Delta^*$ ». La page s'arrête là, le reste du feuillet est
blanc}

\section*{« Th.\ du cycle invariant » : Cas du $H^1$}

\note{titre inscrit au crayon en haut à droite du feuillet de garde
(page~59), qui ne porte rien d'autre, les guillemets étant les siens ; le
feuillet ne reçoit pas de numéro de page. Le tapuscrit qui suit (pages 60 et
61) est donné en entier : les mots tapés qu'il biffe à l'encre sont donnés barrés, ses ajouts
manuscrits comme ajouts de l'éditeur avec une note, les biffures faites à la machine
ne sont pas relevées. Pas de numéro de page tapé sur le feuillet}

\page{60}

Le résultat qui suit, qui résoud le « problème du cycle invariant » dans le cas
du $H^1$, est essentiellement une reformulation d'un résultat de M.~RAYNAUD
[\quad]:
\marginal{une croix au crayon en regard de cette ligne}

\textbf{Théorème} Soient $S$ un trait strictement local, $f\colon X \to S$ un
schéma pro\supplied{jectif} sur $S$, régulier, tel que l'on ait $\underline{O}_S
\xrightarrow{\ \sim\ } f_*(\underline{O}_X)$, (ce qui implique que $X$ est plat
sur $S$, et que $X_0$ est un diviseur sur $X$), $(D_i)_{i \in J}$ la famille
des composantes irréductibles réduites de $X_0$, considérées comme diviseurs
sur $X$, et posons
\note{« jectif » est complété à l'encre au bout de la ligne ; « $(D_i)_{i\in J}$
la famille des » est tapé en interligne, au-dessus d'un « Soient $(D_i)$ »
biffé à la machine, et l'indice $i \in J$ est repris à l'encre ; un trait
vertical à l'encre court dans la marge gauche le long de l'énoncé}
\[
\begin{array}{ll}
(1) & X_0 = \sum_{i \in J} d_i D_i, \qquad \text{où } d_i \geqslant 1 , \\
(2) & d = \operatorname{pgcd}(d_i)_{i \in J}.
\end{array}
\]
Soient
\[
(3) \qquad \Delta \simeq \underline{\mathbf{Z}}^J
\]
la catégorie des diviseurs $D$ sur $X$ tels que $\operatorname{supp} D \subset
X_0$, \struck{et}
\[
(4) \qquad \Gamma \simeq {}/\underline{\mathbf{Z}}.X_0 \hookrightarrow
\operatorname{Pic}(X)
\]
le \struck{est} sous-groupe de $\operatorname{Pic}(X)$ image de $\Delta$ par
l'homomorphisme $D \mapsto \operatorname{cl}(\underline{O}_X(D))$ de $\Delta$ dans
$\operatorname{Pic}(X)$, et considérons l'homomorphisme \struck{composé}
\note{dans (4), le numérateur du quotient est laissé en blanc par la machine
(sans doute $\Delta$, à encrer) ; les $\Delta$, $\Gamma$ et les flèches sont
encrés à la main dans tout le feuillet}
\[
(5) \qquad w\colon \Gamma \longrightarrow \mathrm{NS}(X_0)
\overset{\mathrm{dfn}}{=} \operatorname{Im}(\operatorname{Pic}(X_0) \to
\underline{\mathrm{NS}}_{X_0/k}(k)) ,
\]
composé des homomorphismes évidents
\[
\Gamma \overset{i}{\hookrightarrow} \operatorname{Pic}(X) \longrightarrow
\operatorname{Pic}(X_0) \longrightarrow \mathrm{NS}(X_0) .
\]
Soit enfin $d^t$ le pgcd des degrés des $0$-cycles sur $X_\eta$. Avec ces
notations, on a ce qui suit:
\note{« dfn » est ajouté à l'encre au-dessus du signe $=$ de (5) ; l'indice
$\eta$ de $X_\eta$ est repris à l'encre sur un $y$ tapé}

a) L'entier $d^t$ est aussi égal au pgcd des $\mu_i d_i$, où pour tout $i \in
J$, $\mu_i$ désigne la multiplicité radicielle de $D_i$ sur $k$ (EGA~IV~4.\quad);
en particulier \struck{$d$ divise $d^t$, donc est égal à $1$ si $d^t$ l'est}
\struck{et lui est égal si $k$ est parfait} \struck{\ill{} vrai si $k$ est
parfait} \add{on a $d^t = p^h d$, où $p^h = \operatorname{pgcd}(\mu_i)$},
\marginal{donc $h = 0$ i.e.\ $d = d_t$ si $k$ est parfait}
\note{la fin de la phrase tapée est biffée à l'encre, avec deux corrections
manuscrites elles-mêmes biffées (« et lui est égal si $k$ est parfait »
en interligne, « \ill{} vrai si $k$ est parfait » sous la ligne), puis
remplacée par l'ajout manuscrit « on a $d^t = p^h d$, où $p^h =
\operatorname{pgcd}(\mu_i)$ » ; la marginale, à gauche dans un trait, est reliée
par une flèche à cet ajout. Une croix au crayon en marge de la ligne de
$\mu_i$}

b) L'homomorphisme $w$ est un homomorphisme de $\underline{\mathbf{Z}}$-modules
\emph{de type fini dont le noyau est contenu dans le sous-groupe de torsion de
$\Gamma$}. Ce dernier est isomorphe canoniquement à
$\underline{\mathbf{Z}}/d\underline{\mathbf{Z}}$, avec comme générateur l'image dans
$\Gamma$ du cycle $(1/d)X_0 = \sum_{i \in J} (d_i/d) D_i$. Donc
$\operatorname{Ker} w$ est cyclique
\note{le passage en italique est souligné à l'encre ; un trait vertical court
dans la marge gauche le long de b). Le feuillet s'arrête ici ; la phrase
continue page 61}

\end{document}
